\documentclass[11pt]{article}
\usepackage[T1]{fontenc}
\usepackage[utf8]{inputenc}
\usepackage{lmodern}
\usepackage[margin=0.9in]{geometry}
\usepackage{microtype}
\usepackage{amsmath,amssymb,amsthm}
\usepackage{graphicx}
\usepackage[hidelinks]{hyperref}

\newtheorem{theorem}{Theorem}
\newtheorem{remark}{Remark}
\newcommand{\D}{\mathcal D}
\newcommand{\RZ}{\mathcal{RZ}}

\title{A Resonant Counterexample to Closedness of $\boldsymbol{\pm iA+B}$}
\author{Solution packet for arXiv:2601.19075}
\date{9 August 2026}

\begin{document}
\maketitle

\noindent\textbf{Status.} Candidate counterexample, likely valid; a full
negative answer to the closedness uncertainty in Remark 6.3(iii) of the source
paper.

\section{Source question and setting}
Roidos studies, on $Y=L^p(0,T;X)$, the pointwise extension of a closed operator
$A$ on $X$ together with
\[
  Bu=u',\qquad \D(B)=W^{1,p}_0(0,T;X),
\]
where $u(0)=0$ \cite{Roidos}.  Remark 6.3(iii), on PDF page 29, says that it is
not known whether $\pm iA+B$ is closed.  The same remark notes that positive
definite self-adjoint operators on Hilbert space satisfy the relevant
sectorial and $R$-strip-type hypotheses.

\begin{figure}[h]
  \centering
  \includegraphics[width=0.96\textwidth]{figures/open_problem_crop.png}
  \caption{Remark 6.3(iii) states the closedness uncertainty; part (v) records
  the admissible Hilbert/self-adjoint setting.}
\end{figure}

\section{Proof intuition}
The two unbounded terms can resonate.  On an eigenvector with $Ae_n=ne_n$,
the phase $e^{-\sigma int}$ makes the large terms in
$u'+\sigma iAu$ cancel exactly.  We superpose these resonant modes with
coefficients $1/n$.  The forcing and the solution remain $\ell_2$-valued, but
applying $A$ removes the factor $1/n$ and produces the non-square-summable
constant sequence.  Finite-mode truncations are in the classical intersection
domain, so their graph limit proves that the unclosed sum genuinely has a
larger closure domain.

\section{Counterexample}
\begin{theorem}
Let $1<p<\infty$, $T>0$, $X=\ell_2(\mathbb N)$, and define
\[
  \D(A)=\left\{x=(x_n):\sum_{n=1}^{\infty}n^2|x_n|^2<\infty\right\},
  \qquad (Ax)_n=nx_n.
\]
Extend $A$ pointwise to $Y=L^p(0,T;X)$, and let
$B u=u'$ with $\D(B)=W^{1,p}_0(0,T;X)$.  For each
$\sigma\in\{-1,1\}$, the operator
\[
  S_\sigma=B+\sigma iA,
  \qquad
  \D(S_\sigma)=W^{1,p}_0(0,T;X)\cap L^p(0,T;\D(A)),
\]
is not closed on $Y$.  Nevertheless, $A$ is positive definite and
self-adjoint, $A\in\mathcal P(0)$, and $A\in\RZ_c$ for every $c>0$ in the
notation of the source.
\end{theorem}

\begin{proof}
Let $a_n=1/n$.  For $t\in[0,T]$ define
\[
  u_n(t)=t a_n e^{-\sigma int},
  \qquad
  f_n(t)=a_n e^{-\sigma int}.
\]
Since $a=(a_n)\in\ell_2$, the functions $u$ and $f$ belong to $Y$; indeed,
\[
  \|u(t)\|_{\ell_2}=t\|a\|_{\ell_2},
  \qquad
  \|f(t)\|_{\ell_2}=\|a\|_{\ell_2}.
\]

Let $P_N$ be the projection onto the first $N$ coordinates and put
$u^{(N)}=P_Nu$, $f^{(N)}=P_Nf$.  These finite-coordinate functions belong to
$\D(S_\sigma)$ and vanish at $t=0$.  Coordinatewise differentiation gives
\[
  \frac{d}{dt}\bigl(t a_ne^{-\sigma int}\bigr)
  +\sigma in\,t a_ne^{-\sigma int}
  =a_ne^{-\sigma int},
\]
and hence
\[
  S_\sigma u^{(N)}=f^{(N)}.
\]

The truncations converge in $Y$.  More explicitly,
\[
 \|u-u^{(N)}\|_Y
 =\left(\frac{T^{p+1}}{p+1}\right)^{1/p}
   \left(\sum_{n>N}\frac1{n^2}\right)^{1/2}\longrightarrow0,
\]
while
\[
 \|f-f^{(N)}\|_Y
 =T^{1/p}\left(\sum_{n>N}\frac1{n^2}\right)^{1/2}
 \longrightarrow0.
\]
Thus $(u^{(N)},S_\sigma u^{(N)})\to(u,f)$ in $Y\times Y$.

However, for every $t>0$,
\[
  (Au(t))_n=nt a_ne^{-\sigma int}=t e^{-\sigma int},
\]
which is not an element of $\ell_2$.  Therefore
$u\notin L^p(0,T;\D(A))$ and in particular $u\notin\D(S_\sigma)$.  The graph
of $S_\sigma$ is not closed.

It remains only to check the source hypotheses.  The diagonal operator $A$ is
positive definite and self-adjoint.  For $\lambda\ge0$,
\[
  \|(A+\lambda)^{-1}\|=\sup_{n\ge1}\frac1{n+\lambda}
  =\frac1{1+\lambda},
\]
so $A\in\mathcal P(0)$.  If $|\operatorname{Im}z|\ge c$, then
\[
  \|(A+z)^{-1}\|=\sup_{n\ge1}\frac1{|n+z|}\le\frac1c.
\]
On a Hilbert space uniform boundedness of an operator family implies
$R$-boundedness (apply orthogonality of the Rademacher functions).  Hence the
resolvent family required in the definition of $\RZ_c$ is $R$-bounded, and
$A\in\RZ_c$ for every $c>0$.
\end{proof}

\section{What the counterexample settles}
The answer to the literal closedness uncertainty is negative, and it remains
negative for both signs, every $1<p<\infty$, every $T>0$, and a positive
self-adjoint operator on a Hilbert space.  In particular, the use of
$\overline{\pm iA+B}$ in Section 5 of the source is not merely technical under
the stated operator assumptions.

The result does not identify $\D(\overline{\pm iA+B})$ in Sobolev terms and
does not describe the spaces $E_2^\pm$.  It also does not claim that the extra
embedding assumption in Theorem 6.2 holds for this example.

\section{Verification and novelty check}
The argument is exact and non-computational.  The critical checks are:
(1) the sign in the resonant exponential, (2) convergence in the ambient
$L^p$ space rather than the intersection graph norm, and (3) failure of the
pointwise $\D(A)$ condition for every $t>0$.

A bounded search on 9 August 2026 used the exact title and arXiv id, the quoted
sentence from Remark 6.3(iii), and combinations of
\texttt{closed sum}, \texttt{derivation operator}, \texttt{strip-type}, and
\texttt{$\pm iA+B$}.  It found no later paper or exact answer to this May 2026
uncertainty.  Novelty is therefore provisional pending expert review of the
broader operator-sum literature.

\section{Human-review recommendation}
Promote after checking that Remark 6.3(iii) is intended under the Section 5
hypotheses exactly as written.  The mathematical graph-limit construction is
elementary and appears robust; the principal review risk is scope and novelty,
not the cancellation calculation.

\begin{thebibliography}{9}
\bibitem{Roidos}
N. Roidos,
\emph{Strip-type operators and abstract Cauchy problems},
\href{https://arxiv.org/abs/2601.19075}{arXiv:2601.19075}, current version dated
26 May 2026.
\end{thebibliography}

\end{document}
